\chapter{Theoretical Background}
\section{CNNs and transfer learning}
A CNN slides small learned filters over an image to detect local patterns such as edges and spots. Stacking many such layers lets the network recognise larger shapes and textures. EfficientNet-B0 is built from efficient blocks called MBConv, which include a small ``attention'' step (squeeze-and-excitation) that tells the network which feature channels matter most. EfficientNet scales depth $d$, width $w$, and image resolution $r$ together as $d=\alpha^\phi$, $w=\beta^\phi$, $r=\gamma^\phi$ with $\alpha\beta^2\gamma^2\approx2$~\cite{tan2019efficientnet}. We keep the B0 size.

On top of the network we place a small \emph{classifier head}: dropout, an optional hidden layer, and ten outputs, one per class. The softmax function $p_c=\exp(z_c)/\sum_j\exp(z_j)$ turns the ten outputs $z_c$ into probabilities that add up to one. Layers that are ``frozen'' keep their pretrained values; ``unfrozen'' layers are retrained on tomato leaves. Retraining later layers helps because they hold the most task-specific features~\cite{yosinski2014transfer}.

\section{Training and regularisation}
Training minimises the cross-entropy loss $\mathcal L=-N^{-1}\sum_i\sum_c q_{ic}\log p_{ic}$, which is small when the model gives high probability to the correct class. Several settings help the model generalise instead of memorising:
\begin{itemize}
\item \textbf{Dropout} randomly switches off some connections during training.
\item \textbf{Weight decay} keeps the network's weights small.
\item \textbf{Label smoothing} softens the target answers slightly, $q_{ic}=(1-\epsilon)y_{ic}+\epsilon/C$, so the model is less over-confident.
\item \textbf{Augmentation} shows randomly flipped, cropped, and colour-changed versions of each image.
\end{itemize}
The baseline uses the Adam optimiser; the searched recipes use AdamW, a version of Adam with built-in weight decay. The learning rate rises briefly at the start (warm-up) and then decreases smoothly (cosine schedule). Training stops early when validation macro-F1 stops improving, and the best checkpoint is kept.

\section{Classification metrics}
For class $c$, $TP_c$ counts correct detections, $FP_c$ false alarms, and $FN_c$ missed cases:
\begin{align}
P_c&=\frac{TP_c}{TP_c+FP_c}, & R_c&=\frac{TP_c}{TP_c+FN_c}, & F_{1,c}&=\frac{2P_cR_c}{P_c+R_c},\\
F_{1,\mathrm{macro}}&=\frac{1}{C}\sum_cF_{1,c}, & \mathrm{Accuracy}&=\frac{1}{N}\sum_i\mathbf1[\hat y_i=y_i].
\end{align}
Precision $P_c$ asks ``when the model says class $c$, how often is it right?''; recall $R_c$ asks ``how many real class-$c$ leaves did it find?''; F1 balances the two. Macro-F1 averages F1 over all classes, so every disease counts equally. The ROC curve shows how well the model separates one class from the rest at every confidence threshold, and its area (AUC) is 100\% for perfect separation.

\section{Statistical checks}
\emph{Bootstrap confidence intervals} give a plausible range for a score by repeatedly re-sampling the test images~\cite{efron1979bootstrap}. The \emph{McNemar test} compares two models on the same test images. It looks only at the images where exactly one model is right and checks whether this split could happen by chance~\cite{dietterich1998tests}. A very small $p$-value means the difference is real.

\section{Calibration, explanations, and robustness}
A model is well \emph{calibrated} when its confidence matches how often it is right. Expected calibration error (ECE) measures the gap, averaged over confidence bins $B_m$~\cite{guo2017calibration}:
\begin{equation}
\mathrm{ECE}=\sum_{m=1}^{M}\frac{|B_m|}{N}\left|\mathrm{acc}(B_m)-\mathrm{conf}(B_m)\right|.
\end{equation}
\emph{Grad-CAM} produces a heat map showing which parts of the image most influenced the prediction~\cite{selvaraju2017gradcam}. \emph{Corruption tests} add blur, noise, or darkness to test images to see how the model copes with poorer photos~\cite{hendrycks2019corruptions}.
